\section{Syntax of a Minimalistic Smart Contract Language}\label{sec:syntax}

The pseudocode used in Sec.~\ref{sec:examples} is human-readable,
but still too complicated for analysis. In particular, we want to simplify
the definition of the abstract analyses as much as possible. Therefore, we
assume to analyze programs written in a minimalistic language where,
for instance, variables are only assigned once and commands operates on variables only.
These constraints are not restrictive, but make the code under analysis longer
and require to introduce temporary variables. This is why this minimalistic language
has not been used in Sec.~\ref{sec:examples}. The translation of the smart contracts
from Sec.~\ref{sec:examples} in this language can be found in \textbf{[cite Github repository eventually]}.

\begin{definition}[Expressions and Commands]\label{def:exp_com}
  \emph{Expressions} and \emph{commands} are defined by the following grammar,
  where $c$ is a literal constant of the language, $v$ is a variable, $\mathit{aop}$ is an arithmetic
  operator and $\mathit{cop}$ is a comparison operator:
  \begin{align*}
    \exp::=&\ \mathit{c}\mid\mathit{v}\mid\field{v}{f}\mid\aop{v}{v}\mid\readmap{v}{v}{v}\\
    &\mid\new{C}{v^*}\mid\size{v}\mid\last{v}\mid\random{v}\\
    \com::=&\ \assign{\mathit{v}}{\exp}\mid\assign{\field{v}{f}}{v}\mid\writemap{v}{v}{v}
    \mid\require{v}{v}\\
    &\mid\push{v}{v}\mid\receive{v}{v}\mid\call{v}{m}{v^*}\\
    &\mid\ife{v}{v}{\com}{\com}\mid\for{\mathit{v}}{v}{\com}\mid\sequence{\com}{\com}\mid\nop.
  \end{align*}
  %
  The set of all expressions is $\Exp$ and the set of all commands is $\Com$.
\end{definition}

\begin{definition}[Program]\label{def:program}
  A \emph{program} is a collection of \emph{classes}.
  Each class $C$ has a superclass, some \emph{field signatures} $C.f$
  and some \emph{methods} $C.m$, declared as
  \[
  C.m(p_1:t_1,\ldots,p_n:t_n)\ \<let>\ l_1:t_1',\ldots,l_q:t_q'\ \<in>\ \com,
  \]
  where $\com$ is the \emph{body} of the method, $p_1\cdots p_n$ are its
  formal parameters, of type $t_1\cdots t_n$, and $l_1\cdots l_q$ are its
  local variables, of type $t_1'\cdots t_q'$. Local variables include the
  implicit variables \<this> (the \emph{receiver} of $C.m$),
  \<caller> (the contract that called $C.m$) and \<now> (the current time).
  If $C.m$ is \<payable>, local variables include also
  \<amount> (the cryptocurrency paid for calling $C.m$).
  A class that extends \<Contract>
  (possibly indirectly) is called a \emph{smart contract}.
  A \emph{constructor} is a method $C.C$ named like the class $C$ that declares it.
  A type is either \emph{primitive} (such as \<int>)
  or a class name $C$. Generic class types (such as lists) are written as
  \<List<t\texttt{>}>, where \<t> is the type of the elements of the list.
  %The variables in scope inside a method
  %form its \emph{typing}\ie a map from each variable in scope to
  %its declared type. Namely, the typing $\tau_{C.m}$ of the method $C.m$ above is the map
  %$\tau_{C.m}=\{p_1\mapsto t_1,\ldots,p_n\mapsto t_n,l_1\mapsto t_1',\ldots,l_q\mapsto t_q'\}$.
  Fields $C.f$ can be \<final>, if they are assigned exactly once in each
  constructor of $C$.
  We assume that programs are type-checked and in static single assignment form\ie local
  variables are assigned exactly once along each execution path.
\end{definition}
%
%Fields and methods can have modifiers such as \<public>, \<private>,
%\<final> and \<payable>, that we do not introduce formally and have the
%usual meaning as, for instance, in Java or Solidity. For simplicity, we
%assume that methods do not return any value.
%We write $\tau_{C.m}(\exp)$ for the static type of an expression $\exp$ occurring inside
%method $C.m$. Its definition is not formalized here, since it is the usual recursive
%definition used by any strongly-typed language.
